\section{Esercizio A -- Sistema LTI a tempo continuo}

Un sistema dinamico $\Sigma=(T,X,U,Y,\mathcal{U},\mathcal{Y},\phi,\eta)$ è una
collezione di insiemi, spazi e funzioni: $T$ è l'insieme dei tempi, $X$ lo spazio
di stato, $U$ e $Y$ le immagini di ingresso e uscita, $\mathcal{U}$ e $\mathcal{Y}$
gli insiemi delle funzioni di ingresso e di uscita ammissibili, $\phi$ la funzione
di transizione di stato (forma esplicita del sistema) ed $\eta$ la mappa di uscita.

Il sistema in esame è SISO (Single-Input Single-Output): ingresso e uscita sono
scalari. È inoltre \emph{proprio}, poiché la matrice $D$ è identicamente nulla e
quindi l'ingresso non influenza istantaneamente l'uscita. Le matrici sono:
\[
A=\begin{bmatrix}
-\tfrac{179}{4} & -\tfrac{331}{4} & \tfrac{205}{8} & -\tfrac{409}{8}\\[2pt]
\tfrac{183}{4} & \tfrac{327}{4} & -\tfrac{213}{8} & \tfrac{409}{8}\\[2pt]
-\tfrac{179}{2} & -\tfrac{331}{2} & \tfrac{205}{4} & -\tfrac{413}{4}\\[2pt]
-\tfrac{179}{2} & -\tfrac{327}{2} & \tfrac{209}{4} & -\tfrac{409}{4}
\end{bmatrix},\quad
B=\begin{bmatrix}-\tfrac12\\[2pt]\tfrac12\\[2pt]-1\\[2pt]-1\end{bmatrix},\quad
C=\begin{bmatrix}-10 & 0 & 0 & 5\end{bmatrix}.
\]

\subsection{Modi naturali}
I modi naturali sono funzioni reali di variabile reale che descrivono l'andamento
della risposta libera; si presentano nella forma $e^{\lambda_i t}$, con $\lambda_i$
l'$i$-esimo autovalore della matrice dinamica $A$. Gli autovalori sono
\[
\lambda_1=-7,\qquad \lambda_2=-3,\qquad \lambda_{3,4}=-2\pm\tfrac{j}{2}.
\]
Due reali e distinti, due complessi coniugati. Per ogni autovalore la molteplicità
algebrica coincide con la geometrica, quindi $A$ è diagonalizzabile. I modi associati
agli autovalori reali sono $e^{-7t}$ ed $e^{-3t}$; quelli legati alla coppia complessa
hanno forma $e^{\sigma t}\cos(\omega t)$ ed $e^{\sigma t}\sin(\omega t)$ con
$\sigma=\operatorname{Re}[\lambda]$ e $\omega=\operatorname{Im}[\lambda]$, cioè
$e^{-2t}\cos\!\big(\tfrac{t}{2}\big)$ ed $e^{-2t}\sin\!\big(-\tfrac{t}{2}\big)$.

\mmacap{img-003}{Andamento del 1\textordmasculine\ modo naturale}
\mmacap{img-004}{Andamento del 2\textordmasculine\ modo naturale}
\mmacap{img-005}{Andamento del 3\textordmasculine\ modo naturale}
\mmacap{img-006}{Andamento del 4\textordmasculine\ modo naturale}

\subsection{Risposta libera}
La risposta libera si ottiene da uno stato iniziale assegnato $x_0$ con ingresso
identicamente nullo, ed è una combinazione lineare dei modi naturali. Sullo stato vale
$x_l(t)=e^{At}x_0$; conviene diagonalizzare $A$ tramite una matrice di cambiamento di
base $T$ tale che $\Lambda=T^{-1}AT$, da cui $x_l(t)=T\,e^{\Lambda t}z_0$ con
$z_0=T^{-1}x_0$.

Per la coppia di autovalori complessi coniugati si costruisce una matrice $T$
``ibrida'', usando parte reale e immaginaria di un autovettore: in questo modo $T$
resta reale e $\Lambda$ assume la forma canonica reale con un blocco di Jordan
$2\times2$.

\mma[0.5]{img-013}

Si proietta lo stato iniziale $x_0=\begin{bmatrix}-1&0&-3&0\end{bmatrix}^{\!\top}$ sulle
colonne di $T$ ottenendo $z_0=T^{-1}x_0$; quindi la risposta libera sullo stato si
calcola come $x_l(t)=T\,\mathrm{MatrixExp}[\Lambda t]\,z_0$ e quella sull'uscita come
$y_l(t)=C\,x_l(t)$.

\mmacap{img-016}{Risposta libera sull'uscita confrontata con quella sullo stato}

\subsection{Configurazione degli stati iniziali}
Poiché la risposta libera è combinazione lineare dei modi naturali, una scelta
opportuna di $x_0$ può ``accendere'' o ``spegnere'' determinati modi. Scegliendo $x_0$
come combinazione lineare delle sole colonne di $T$ associate ai modi che si vogliono
evidenziare, le restanti componenti di $z_0$ si annullano e quei modi non compaiono.

\mmacap{img-018}{Accensione dei modi legati alle prime due colonne di $T$}
\mmacap{img-020}{Accensione dei modi legati alle ultime due colonne di $T$}
\mmacap[0.48]{img-021}{Solo prima colonna}
\mmacap[0.48]{img-022}{Solo seconda colonna}
\mmacap[0.48]{img-023}{Solo terza colonna}
\mmacap[0.48]{img-024}{Solo quarta colonna}

\subsection{Funzione di trasferimento, poli e zeri}
La funzione di trasferimento è la funzione di variabile complessa $G(s)$ tale che
$Y(s)=G(s)U(s)$, calcolata come $G(s)=C(sI-A)^{-1}B+D$ (qui $D=0$). Normalizzando
numeratore e denominatore rispetto al coefficiente di grado massimo si ottiene un
denominatore
\[
d(s)=s^4+14s^3+\tfrac{261}{4}s^2+\tfrac{253}{2}s+\tfrac{357}{4}.
\]

\mma[0.65]{img-025}

La correttezza si verifica con \texttt{StateSpaceModel}, che restituisce la forma di
Rosenbrock del sistema. I \emph{poli} sono le radici del denominatore (sempre un
sottoinsieme dei modi naturali); qui coincidono con gli autovalori:
\[
s=-7,\quad s=-3,\quad s=-2\pm\tfrac{j}{2}.
\]
Poiché i poli coincidono con gli autovalori, nella risposta forzata compaiono tutti i
modi naturali. Gli \emph{zeri} sono le radici del numeratore.

\subsection{Risposta al gradino unitario}
La risposta forzata al gradino è $Y(s)=G(s)\,U(s)$ con $\Lapl[1(t)]=\tfrac1s$.
Scomponendo $Y(s)$ in fratti semplici, la componente legata al polo $s=0$ è la
\emph{risposta a regime}, mentre quelle legate ai poli del sistema costituiscono la
\emph{risposta transitoria}. I residui si calcolano con la formula di Heaviside
\[
C_{ij}=\frac{1}{(v_i-j)!}\lim_{s\to p_i}\Big[(s-p_i)^{v_i}F(s)\Big].
\]
Per la coppia complessa è sufficiente calcolare $C_4$: $C_5=\overline{C_4}$.

\mmacap[0.85]{img-036}{Risposta forzata al gradino con le componenti $y_{ss}(t)$ e $y_{tr}(t)$}

\subsection{Risposta al segnale periodico elementare}
Si considera $u(t)=A\sin(\omega t+\psi)\,1(t)$ con la terna $A=2,\ \omega=2,\ \psi=0$.
Costruite $U(s)$ e $Y(s)$, si scompone in fratti semplici e, con la formula di
Heaviside, si separano la componente di regime (legata all'ingresso) e quella
transitoria (legata ai modi).

\mmacap[0.7]{img-038}{Risposta forzata al segnale sinusoidale e sue componenti}

Il sistema distorce il segnale in ampiezza e fase. Per quantificare i fattori di
distorsione si usa la forma \emph{amplitude-phase} $X\sin(\omega t+\psi)$, eguagliata
alla componente di regime nell'ipotesi di ingresso $\sin(t)\,1(t)$; dei fratti semplici
interessano solo quelli legati all'ingresso.

\mmacap[0.7]{img-050}{Confronto fra ingresso e componente di regime (forma amplitude-phase)}

Risolvendo si trova $X<1$: il segnale è attenuato in ampiezza e mostra un anticipo di
fase.

\subsection{Modello ARMA equivalente e risposta alla rampa}
Il modello ARMA (Auto-Regressive Moving Average) è una rappresentazione implicita
ingresso/uscita, un'equazione differenziale della forma
\[
y^{(n)}+a_1y^{(n-1)}+\dots+a_ny=b_0u^{(n)}+b_1u^{(n-1)}+\dots+b_nu,
\]
ricavabile da $G(s)=\tfrac{Y(s)}{U(s)}$. Le condizioni iniziali sull'uscita compatibili
con $x_0=\begin{bmatrix}-3&2&-3&3\end{bmatrix}^{\!\top}$ si ottengono dalla matrice di
osservabilità: $[\,y(0)\ \dot y(0)\ \ddot y(0)\ \dddot y(0)\,]^{\top}=\mathrm{Ob}\,x_0$,
con $\mathrm{Ob}=[\,C;\,CA;\,CA^2;\,CA^3\,]$. Per il sistema in esame:
\[
\dddot{\!\!\;y}^{\,(4)}+14\,\dddot y+\tfrac{261}{4}\ddot y+\tfrac{253}{2}\dot y+\tfrac{357}{4}y
=5\,\dot u-5\,u.
\]

\mma[0.85]{img-052}

Antitrasformando e confrontando la risposta libera ottenuta dal modello I/U con quella
ricavata dalla definizione I/S/U si verifica che le due coincidono, a conferma della
correttezza dei calcoli.

\paragraph{Risposta alla rampa unitaria.}
Con $\Lapl[t\,1(t)]=\tfrac{1}{s^2}$, si sostituiscono nel modello ARMA le condizioni
iniziali e la trasformata dell'ingresso, separando la componente transitoria (risposta
libera inclusa) da quella legata algebricamente all'ingresso:
\[
y_{ss}(t)=\frac{17260}{127449}-\frac{20\,t}{357}.
\]

\subsection{Stato iniziale che annulla il transitorio}
Si cerca $x_0$ tale che la risposta al gradino coincida con il suo valore di regime
(assenza di transitorio): occorre che la somma di risposta libera e risposta transitoria
sia identicamente nulla. Imponendo l'annullamento del numeratore (polinomio di terzo
grado nelle quattro incognite $x_i$) si ottiene
\[
x_0=\begin{bmatrix}\tfrac{2}{357}&-\tfrac{2}{357}&\tfrac{4}{357}&0\end{bmatrix}^{\!\top}.
\]
La verifica con \texttt{OutputResponse} applicato alla forma di Rosenbrock e il
confronto con $\mathrm{FdT}[0]$ danno valori coincidenti.

\subsection{Risposta al segnale $u(t)=1(-t)$}
L'ingresso $1(-t)$ vale $1$ per $t<0$ e $0$ per $t\ge0$. Lo si interpreta come un segnale
applicato in un passato remoto ($t\to-\infty$): la componente transitoria si è già
estinta e per $t<0$ resta solo la componente di regime. Per $t\ge0$ l'ingresso è nullo e
il sistema evolve liberamente a partire dalle condizioni all'istante $t=0$, ricavate da
$x_0=\mathrm{Ob}^{-1}\cdot\text{(condizioni iniziali)}$. La risposta complessiva si
descrive con un costrutto \texttt{Piecewise}.

\mmacap[0.55]{img-080}{Risposta del sistema all'ingresso $u(t)=1(-t)$}
